\begin{lexlean}{RelativePosition}
\useglossary{lexlean.std.bool@1.2.0}
\useglossary{lexlean.std.nat@1.2.0}
\title{Natural number}
\begin{semanticmodule}
definition{
  name: sourceRow,
  parameters: [positions : Nat, row : Nat, plane : Nat],
  result: Nat,
  body: ((prim subtract (prim subtract positions row : Nat) 1 : Nat) + plane),
  executable: true
}

definition{
  name: sourceWord,
  parameters: [index : Nat, width : Nat, positions : Nat],
  result: Nat,
  body: ((prim multiply (sourceRow positions (prim remainder (prim quotient index width 0 : Nat) positions 0 : Nat) (prim quotient (prim quotient index width 0 : Nat) positions 0 : Nat)) width : Nat) + (prim remainder index width 0 : Nat)),
  executable: true
}

theorem{
  name: everySourceRowIsBounded,
  parameters: [positions : Nat, row : Nat, plane : Nat],
  statement: ((row < positions) -> ((plane < positions) -> ((sourceRow positions row plane) < (prim subtract (positions + positions) 1 : Nat)))),
  proof: linear_arithmetic{definitions: [sourceRow]},
  axioms: [Quot.sound, propext]
}

theorem{
  name: firstRelativeWord,
  parameters: [],
  statement: ((sourceWord 0 3 7) = 18),
  proof: reflexivity{}
}

theorem{
  name: reverseRelativeRow,
  parameters: [],
  statement: ((sourceWord 20 3 7) = 2),
  proof: reflexivity{}
}

theorem{
  name: lastRelativeWord,
  parameters: [],
  statement: ((sourceWord 146 3 7) = 20),
  proof: reflexivity{}
}
\end{semanticmodule}
\end{lexlean}
